% 核心结论：同轴长直螺线管的交链面积和环形磁芯的闭合磁路
% 图原型：纯几何／对比。物理网络按本题节点、端口和符号构建。
\documentclass[tikz,border=7pt]{standalone}
\usepackage{fontspec,xeCJK,amsmath,amssymb}
\IfFontExistsTF{Times New Roman}{\setmainfont{Times New Roman}}{\setmainfont{TeX Gyre Termes}}
\IfFontExistsTF{Heiti SC}{\setCJKmainfont[AutoFakeBold=1.5]{Heiti SC}}{\setCJKmainfont[AutoFakeBold=1.5]{FandolHei-Regular}}
\usetikzlibrary{arrows.meta,calc,angles,quotes}
\definecolor{ink}{HTML}{243745}
\definecolor{blue}{HTML}{275DAB}
\definecolor{green}{HTML}{177C63}
\definecolor{red}{HTML}{BD4935}
\definecolor{purple}{HTML}{715393}
\begin{document}
\begin{tikzpicture}[font=\fontsize{9.5}{12}\selectfont,text=ink,
 v/.style={line width=1.1pt,-{Stealth[length=2.4mm,width=1.5mm]}},
 axis/.style={ink!65,line width=.45pt,-{Stealth[length=1.7mm,width=1.1mm]}},
 aux/.style={ink!40,densely dashed,line width=.5pt},
 lab/.style={fill=white,inner sep=2pt},
 card/.style={draw=blue!25,fill=blue!3,rounded corners=3pt,text width=6cm,align=left,inner sep=10pt}]
\path[use as bounding box] (-8,-3.35) rectangle (8,3.65);
\node[font=\fontsize{10.5}{13}\selectfont\bfseries,text width=15cm,align=center] at (0.00000,3.22000) {M6-H21　同轴长直螺线管的交链面积和环形磁芯的闭合磁路};
\begin{scope}[shift={(-4,0)}]
\draw[ink,fill=blue!5] (0,0) circle(1.75);
\draw[blue,fill=green!12] (0,0) circle(1);
\node[] at (0.00000,0.00000) {$S_1$};
\node[] at (1.35000,0.00000) {$S_2$};
\node[] at (0.00000,2.25000) {同轴长螺线管的横截面};
\node[] at (0.00000,-2.40000) {$M=\mu N_1N_2S_1/l$，$S_1\le S_2$};
\end{scope}
\begin{scope}[shift={(4,0)}]
\fill[blue!6,even odd rule] (0,0) circle(1.85) (0,0) circle(.85);
\draw[ink,thick] (0,0) circle(1.85);
\draw[ink,thick] (0,0) circle(.85);
\draw[axis] (0.00000,0.00000) -- (1.85000,0.00000);
\node[] at (0.40000,0.25000) {$a$};
\node[] at (2.10000,0.55000) {$b$};
\draw[blue,densely dashed] (0,0) circle(1.3);
\node[text=blue] at (-0.60000,1.50000) {$r$};
\draw[green,thick] (-1.95000,-0.40000) -- (-1.25000,-0.20000);
\draw[green,thick] (-1.95000,-0.10000) -- (-1.25000,0.10000);
\draw[green,thick] (-1.95000,0.20000) -- (-1.25000,0.40000);
\draw[green,thick] (1.25000,-0.40000) -- (1.95000,-0.20000);
\draw[green,thick] (1.25000,-0.10000) -- (1.95000,0.10000);
\draw[green,thick] (1.25000,0.20000) -- (1.95000,0.40000);
\node[] at (0.00000,-2.40000) {矩形截面高度 $h$；磁场 $B(r)\propto1/r$。};
\end{scope}
\node[text width=15cm,align=center] at (0.00000,-3.00000) {绕组端口、电流与电压方向按图示约定；磁场模型的近似条件见正文。};

\end{tikzpicture}\end{document}